\documentclass[10pt,a4paper]{amsart}
\usepackage[margin=19mm]{geometry}
\usepackage{amsmath,amssymb,mathtools}
\usepackage[scr=rsfs,cal=euler]{mathalfa}
\usepackage{xfrac}
\usepackage[unicode,hidelinks,bookmarks=false]{hyperref}
\newcommand{\ie}{i.e.\ }
\newcommand{\cf}{cf.\ }
\newcommand{\resp}{respectively\ }
\newcommand{\Subsection}{\subsection}
\providecommand{\nobreakdash}{\nobreak\hbox{-}}
\setlength{\emergencystretch}{2em}
\multlinegap0pt
\usepackage{bm}
\usepackage{bbm}
\usepackage{dsfont}
\usepackage{xspace}
\usepackage{cancel}
\usepackage{mathtools}
\usepackage{amsthm}
\makeatletter
\@ifundefined{thm}{\newtheorem{thm}{Thm}}{}
\@ifundefined{proposition}{\newtheorem{proposition}[thm]{Proposition}}{}
\makeatother
\makeatletter
\expandafter\ifx\csname varthm\endcsname\relax
\newenvironment{varthm}[1]{\def\varthmname{#1}\begin{thevarthm}}{\end{thevarthm}\def\varthmname{}}
\fi
\makeatother
\title[On homological properties of conic-line arrangements with si]{On homological properties of conic-line arrangements with simple singularities}
\author{Artur Bromboszcz}
\date{Source: arXiv:2606.31452v1 (CC BY 4.0)}
\begin{document}
\maketitle

\noindent\emph{Source and licence.} On homological properties of conic-line arrangements with simple singularities. Version arXiv:2606.31452v1 (CC BY 4.0), distributed under the Creative Commons Attribution 4.0 International licence, as recorded in the arXiv licence metadata for this exact version and shown on the abstract page https://arxiv.org/abs/2606.31452v1. Full rights notice: licenses/arxiv-2606.31452-CC-BY-4.0.txt.\par\smallskip

\noindent\textit{Editorial scope.} Counted unit: the Proposition whose source label is \texttt{391} in the article (source0.tex lines 539--546, source label \texttt{391}), delivered together with the article's own complete original proof of it (source0.tex lines 548--650, one \texttt{proof} environment of 103 lines). No mathematics is written, repaired or re-derived: statement and proof are the article's own text line for line. Required context retained uncounted: none. Every definition, display and label of those lines is retained unchanged. Omitted from this package and not redistributed: 1--538, 547--547, 651--2395. Nothing inside a retained excerpt is omitted or altered: every retained body is delivered exactly as the article's lines are written, so no omission receipt of any kind accompanies this package. Citations: the retained excerpts carry no citation command at all, so no bibliography entry of the article is transcribed and no reference is redistributed. Reference adaptation: none was needed and none was made; every reference inside the delivered unit resolves inside it, so no unresolved reference or citation is produced. Printed slips kept literal: no printed word, symbol, hypothesis or number of the counted statement or of its proof was altered. Layout and class: the article's own class is \texttt{article} with packages including graphicx, geometry, float, empheq, pdflscape, multirow, circledsteps, enumitem, comment, hyperref, amsfonts, amssymb, amsthm, amsmath; the wrapper is the lane's pinned \texttt{amsart} editorial wrapper at 10pt on A4, which restates the theorem-like environments the retained text needs and adds the article's own macro definitions that the retained text needs (none), with extra wrapper packages (bm, bbm, dsfont, xspace, cancel, mathtools). The delivered numbering is the unit's own. Assets: no retained span contains a figure environment, a graphics-inclusion command, a PDF-page inclusion, a table or a TikZ picture; no image file is copied into this package, the package declares no asset dependency and no image file is redistributed with it. Licence: version arXiv:2606.31452v1 of this article is distributed under the Creative Commons Attribution 4.0 International licence, as recorded in the arXiv licence metadata for this exact version and shown on the abstract page https://arxiv.org/abs/2606.31452v1. A full rights notice is delivered at licenses/arxiv-2606.31452-CC-BY-4.0.txt. Characters: every retained excerpt is pure ASCII, so the delivered engine, XeLaTeX with its default Latin Modern text font, reports no missing character.
\begin{proposition}
\label{391}
The weak combinatorics
\[
        (n_2,t_3,n_3)=(3,9,1)
\]
is geometrically realizable over \(\mathbb C\) by an arrangement of four smooth conics.
\end{proposition}

\begin{proof}
Let \(\imath^2=-1\). We consider the following four conics in
\(\mathbb P^2_{\mathbb C}\):
\[
\begin{aligned}
Q_0&: q_0=y^2-4xz=0,\\
Q_1&: q_1=x^2-12\imath xy-146xz+12\imath yz+z^2=0,\\
Q_2&: q_2=256x^2-48\imath xy-137xz+12\imath yz+16z^2=0,\\
Q_3&: q_3=729x^2-108\imath xy-178xz+12\imath yz+9z^2=0.
\end{aligned}
\]
We denote by
\[
        \mathcal C=Q_0\cup Q_1\cup Q_2\cup Q_3
\]
the corresponding conic arrangement.

\noindent
We now describe the singular locus of \(\mathcal C\). The conic \(Q_0\)
has the parametrization
\[
        \phi(t)=(t^2:2t:1).
\]
Substituting this parametrization into the remaining three equations gives
\[
\begin{aligned}
q_1(\phi(t))&=(t^2-12\imath t-1)^2,\\
q_2(\phi(t))&=(16t^2-3\imath t-4)^2,\\
q_3(\phi(t))&=(27t^2-4\imath t-3)^2.
\end{aligned}
\]
The three quadratic factors are square-free and pairwise coprime. Therefore
\(Q_0\) is tangent to each of \(Q_1,Q_2,Q_3\) at two distinct points, and these
six tangency points are pairwise distinct. Consequently, the intersections
\[
        Q_0\cap Q_1,\qquad Q_0\cap Q_2,\qquad Q_0\cap Q_3
\]
contribute six tacnodes.

Next, the conics \(Q_1,Q_2,Q_3\) pass through the point
\[
        P=(0:1:0),
\]
whereas \(P\notin Q_0\), since \(q_0(P)=1\). In the affine chart \(y=1\),
with local coordinates \(u=x/y\) and \(v=z/y\), the linear parts of
\(q_1,q_2,q_3\) at \(P\) are proportional to
\[
        v-u,\qquad v-4u,\qquad v-9u,
\]
respectively. These three tangent directions are distinct, hence \(P\) is an
ordinary triple point of the arrangement.

The remaining pairwise intersections among \(Q_1,Q_2,Q_3\) are the following:
\[
\renewcommand{\arraystretch}{1.4}
\begin{array}{c|c|c}
\text{pair} & \text{node} & \text{tacnode}\\
\hline
(Q_1,Q_2)
&
\left(1:-\dfrac{929\imath}{60}:\dfrac{21}{5}\right)
&
\left(1:-\dfrac{33\imath}{4}:-2\right)
\\
(Q_1,Q_3)
&
\left(1:-\dfrac{151\imath}{12}:10\right)
&
\left(1:-\dfrac{28\imath}{3}:-3\right)
\\
(Q_2,Q_3)
&
\left(1:-\dfrac{3719\imath}{420}:-\dfrac{17}{7}\right)
&
\left(1:\dfrac{5\imath}{12}:6\right).
\end{array}
\]
A direct substitution shows that each of the above points lies on the indicated
pair of conics and on no other component of \(\mathcal C\). Moreover, the local
intersection multiplicities are
\[
\begin{aligned}
I_P(Q_1,Q_2)&=I_P(Q_1,Q_3)=I_P(Q_2,Q_3)=1,\\
I_{N_{12}}(Q_1,Q_2)&=I_{N_{13}}(Q_1,Q_3)=I_{N_{23}}(Q_2,Q_3)=1,\\
I_{T_{12}}(Q_1,Q_2)&=I_{T_{13}}(Q_1,Q_3)=I_{T_{23}}(Q_2,Q_3)=2.
\end{aligned}
\]
Thus \(N_{12},N_{13},N_{23}\) are nodes, while
\(T_{12},T_{13},T_{23}\) are tacnodes.

For each pair of smooth conics the total intersection multiplicity is equal to
\(4\), by B\'ezout's theorem. Hence the points listed above exhaust all pairwise
intersections among \(Q_1,Q_2,Q_3\). Combining this with the six tacnodes
coming from the intersections of \(Q_0\) with \(Q_1,Q_2,Q_3\), we obtain
\[
        n_2=3,\qquad t_3=6+3=9,\qquad n_3=1.
\]
Therefore the weak combinatorics
\[
        (n_2,t_3,n_3)=(3,9,1)
\]
is geometrically realizable over \(\mathbb C\).
\end{proof}

\end{document}
